\section{Dirbtinio intelekto naudojimo savirefleksija}

Pagal KTU 2026 m. šabloną šis skyrius yra privalomas: jame aprašoma, kokie dirbtinio intelekto (DI) įrankiai naudoti, kuriuose darbo etapuose ir kaip buvo tikrinami jų rezultatai.

\subsection{Naudoti DI įrankiai ir kontekstas}

\lipsum[1]

\subsection{Taikymas projekto kūrimo etapuose}

Apibendrinta DI naudojimo apžvalga pateikiama \ref{tab:di-etapai} lentelėje.

\begin{table}[!ht]
\caption{DI panaudojimas projekto kūrimo etapuose}
\label{tab:di-etapai}
\centering
\begin{tabular}{|p{2.6cm}|p{6.8cm}|p{5.8cm}|} \toprule
\textbf{Etapas} & \textbf{DI veikla} & \textbf{Verifikavimo metodas} \\ \midrule
Analizė & \ldots & \ldots \\ \hline
Projektavimas & \ldots & \ldots \\ \hline
Programavimas & \ldots & \ldots \\ \hline
Testavimas & \ldots & \ldots \\ \hline
Dokumentavimas & \ldots & \ldots \\ \bottomrule
\end{tabular}
\end{table}

\subsection{Stiprybės ir trūkumai}

\lipsum[2]
